% 既存の75題を、読む目的に応じて選ぶための案内。
\subsection*{演習の読む経路}\label{節：演習の読む経路}
  初回からすべての問題を解く必要はない.
  初学者は,読んだ型の例題を自分で解いた後に,次の代表12題を順に選べる.
  一度解いたら本文を閉じ,変形を選ぶ理由と必要な条件を説明してみよう.

  \noindent\textbf{【差と比を追う】}
  \begin{itemize}
    \item \bookproblemref{practice-basic}{10}{標準問10}：等差の更新.
    \item \bookproblemref{practice-basic}{5}{標準問5}：等比の更新.
    \item \bookproblemref{practice-basic}{11}{標準問11}：階差を足し戻す.
    \item \bookproblemref{practice-basic}{3}{標準問3}：倍率を階乗でそろえる.
  \end{itemize}
  \noindent\textbf{【更新が簡単になる量を選ぶ】}
  \begin{itemize}
    \item \bookproblemref{practice-basic}{9}{標準問9}：不動点からの差.
    \item \bookproblemref{practice-basic}{4}{標準問4}：同じ更新をする一次式との差.
    \item \bookproblemref{practice-basic}{8}{標準問8}：対数で指数を取り出す.
    \item \bookproblemref{practice-basic}{1}{標準問1}：逆数を使う.
    \item \bookproblemref{practice-basic}{26}{標準問26}：二つの不動点からの差の比.
  \end{itemize}
  \noindent\textbf{【追う対象を組み替える】}
  \begin{itemize}
    \item \bookproblemref{practice-basic}{6}{標準問6}：隣接3項間を等比型へ分ける.
    \item \bookproblemref{practice-basic}{2}{標準問2}：総和の式をずらして引く.
    \item \bookproblemref{practice-basic}{7}{標準問7}：連立する量の和と差.
  \end{itemize}
  係数や見た目が変わった場合の確認には,標準の
  \bookproblemref{practice-basic}{17}{問17},
  \bookproblemref{practice-basic}{20}{問20},
  \bookproblemref{practice-basic}{23}{問23},
  \bookproblemref{practice-basic}{34}{問34}を使う.
  変形を組み合わせる段階では,応用の
  \bookproblemref{practice-advanced}{1}{問1},
  \bookproblemref{practice-advanced}{7}{問7},
  \bookproblemref{practice-advanced}{9}{問9}へ進める.

  \noindent\textbf{【既習者の比較題】}
  \begin{itemize}
    \item \bookproblemref{practice-basic}{4}{標準問4}と\bookproblemref{practice-basic}{24}{問24}：引く量を探す方法と,階差や正規化を使う方法は,何を消しているか.
    \item \bookproblemref{practice-basic}{19}{標準問19}と\bookproblemref{practice-basic}{42}{問42}：積・正規化・対数で,条件と計算量はどう違うか.
    \item \bookproblemref{practice-basic}{23}{標準問23}と\bookproblemref{practice-basic}{25}{問25}：特性根の重解と,階差を重ねる見方はどう結び付くか.
  \end{itemize}
  別解を読んだら,「どちらが短いか」だけでなく,初項や係数を変えても使えるか,例外がどこで分かるかを比較する.

  \noindent\textbf{【付録と往復する】}
  \begin{itemize}
    \item \bookproblemref{practice-entrance}{3}{発展問3}の総和は,\sectionref*{節：数列全体を一つの対象にまとめると,何が見えるか}の母関数の積と比べられる.
    \item \bookproblemref{practice-entrance}{4}{発展問4}の図形の更新は,複素数による回転と行列による更新から見直せる.
    \item \bookproblemref{practice-entrance}{6}{発展問6}と\bookproblemref{practice-entrance}{9}{問9}では,一般項に限らず,減り方や整数の関係を保つ更新を調べる.
  \end{itemize}
